3.4 Dynamic Tube Stage
3.4.1 Gray-Box Modeling Approach
While the static nonlinearity (Section 3.1) and the parasitic models (Sections 3.2 and 3.3) capture the core components of tube Distortion, a truly authentic simulation must account for the dynamic interactions between the signal envelope and the tube's operating parameters. We propose a "gray-box" dynamic tube stage model, implemented as MODEL 2 in our framework, which modulates the transfer function in real-time based on the player's performance.

The model is defined by a pair of coupled modulation equations:
\begin{align}
d[n] &= d_0 + k_F \cdot e_F[n] - k_S \cdot e_S[n] \tag{12} \\
b[n] &= b_0 + k_B \cdot e_S[n] \tag{13} \\
y[n] &= \tanh\bigl(d[n] \cdot (x[n] + b[n])\bigr) \tag{14}
\end{align}

where $e_F[n]$ and $e_S[n]$ are fast and slow envelope followers, respectively. Parameters $k_F$, $k_S$, and $k_B$ govern the sensitivity of the tube to these dynamic signals.

3.4.2 Dynamic Parameters and Psychoacoustics
The dynamic model introduces several high-level behaviors that are central to the "tube experience":

Touch Sensitivity ($k_F, \tau_F$): The fast modulation $k_F$ increases the drive $d[n]$ instantaneously in response to pick transients. This produces a "bloom" effect where the distortion intensifies during the attack before settling, providing the articulate response highly valued by lead guitarists.
Power Supply Sag ($k_S, \tau_S$): The slow modulation $k_S$ reduces the drive as the average signal level increases. This simulates the voltage drop (sag) in the power supply's B+ rail and the current-limiting effects of the plate and screen resistors. Perceptually, this manifests as a natural, "spongy" compression that smooths out the player's dynamics.
Dynamic Operating Point ($k_B$): Perhaps the most sophisticated feature, $k_B$ allows the bias $b[n]$ to shift over time. As shown in Section 3.1, bias determines the H2/H3 harmonic balance. By making this shift dynamic, the "character" of the distortion evolves as a note sustains, moving from a symmetrical to an asymmetrical state.

3.4.3 Figure Descriptions and Interpretation
Figure 1 — Envelopes and Parameter Modulation. This figure demonstrates the real-time tracking of the model. Panel (a) shows the input signal (a burst) and the two envelope followers. Panel (b) reveals the resulting dynamic drive $d(t)$; note the initial overshoot caused by $k_F$ followed by the gradual sag caused by $k_S$. Panel (c) shows the bias $b(t)$ crawling toward a more negative value during the burst, simulating the dynamic accumulation of charge in the cathode and grid circuits.

Figure 2 — Static vs. Dynamic Comparison. Panel (a) highlights the "sag" effect in the time domain using an RMS envelope. Unlike the static model which has a flat response, the dynamic model exhibits a natural gain reduction over the first 100ms of the note. Panel (b) compares the waveform at the start of a note versus its sustained tail; the operating point shift is visible as the waveform becomes progressively more asymmetrical.

Figure 3 — Dynamic Harmonic Balance. This analysis tracks the second (H2) and third (H3) harmonics over the duration of a sustained note. While H3 remains relatively stable due to the constant input level, H2 grows significantly as the dynamic bias $k_B$ shifts the operating point. This "harmonic bloom" is a hallmark of high-quality tube amplification that is missing in purely static waveshaping models.

3.4.4 Application to Tube Archetypes
By adjusting these dynamic coefficients, we can emulate different tube types within the same mathematical framework. For example, a 12AX7 preset is configured with high $k_F$ for maximum touch sensitivity, whereas a 6V6 power tube emulation utilizes a higher $k_S$ and longer $\tau_S$ to emphasize the characteristic "squish" and sag of a vintage power amplifier.
